\section{Introduction}
\label{sec:introduction}

Releasing LLM-based applications safely is hard because their outputs are probabilistic, prone to ``hallucinations,'' and heavily context-dependent~\cite{dobslaw2025testing}. Standard unit and integration tests, designed for deterministic software, miss the behavioral regressions that matter most in production: persona drift, subtle safety violations, and unsupported evidence claims. This is particularly visible in domains such as marketing automation and conversational analytics, where multi-agent orchestrators compose tool calls, retrieval, and natural-language reasoning in a single user-visible response.

This mismatch creates a persistent release question for engineering teams: \textit{when is a new conversational AI release ready for production?} Beyond functional correctness, teams must enforce guardrails consistently and prevent regressions in reasoning behavior. The current ecosystem still lacks automated, data-driven end-to-end testing frameworks that evaluate qualitative dimensions from a realistic user-perception perspective, which limits reliable continuous deployment in LLM systems. Our focus here is narrow: an automated self-testing loop for longitudinal quality-gated deployment of a specific multi-agent application.

To mitigate this gap, we present an empirical testing framework based on \textit{Automated Self-Testing} (historically known as dogfooding \cite{harrison2006dogfooding}), designed to serve as \textit{Quality Gates} prior to the release of complex LLM systems. By creating a continuous loop where the system is intensively evaluated against a static and dynamic ``question bank,'' our framework assesses the build across five empirically grounded dimensions---Task Success Rate (build stability and response utility), Research Context Preservation (multi-turn context continuity), P95 Latency (user experience), Safety Pass Rate (guardrail adherence), and Evidence Coverage (factuality anchoring)---selected to cover both technical robustness (latency, safety) and user-perceived quality (task success, evidence). Based on these metrics, the framework generates a clear release management decision: PROMOTE/HOLD/ROLLBACK.

We validated our method through a longitudinal empirical study (\textit{case study}) using an internally deployed multi-agent orchestration system (i.e., a conversational AI agent with specific marketing and analytics capabilities) in active development. We analyzed 38 evaluation runs with question banks evolving from 59 to 133 scenarios across more than 20 internal system releases. The framework flagged 2 ROLLBACK-grade builds---both confirmed as severe regressions by the gate ablation analysis (Section~\ref{sec:results})---and supported stable quality evolution over a four-week staging lifecycle. Statistical analysis reveals that evidence coverage was the primary determinant of release rejection, while success rate exhibits a controlled decrease as the question bank grows more challenging.

The main contributions of this paper are:
\begin{enumerate}
    \item \textbf{A formalized and scalable self-testing framework} acting as a \textit{quality gate} (decision checkpoint) for Agentic AI release governance.
    \item \textbf{A longitudinal internal case study}, detailing quality-gated evolution across 38 evaluated runs over 20+ rapid releases in staging, including persona-grounded and multi-turn scenarios.
    \item An \textbf{empirically calibrated set of quality acceptance thresholds} for systemic approval (PROMOTE/HOLD/ROLLBACK), accompanied by illustrative pseudocode and sample test cases that enable community adaptation without exposing proprietary assets.
\end{enumerate}

\paragraph{Paper Structure.} The remainder of this paper is organized as follows: Section \ref{sec:related-work} contextualizes our work within the evolving landscape of LLM testing and continuous delivery. Section \ref{sec:framework} details the automated self-testing framework and quality dimensions. Section \ref{sec:case-study} describes the multi-agent conversational AI case study setup. Section \ref{sec:results} presents the empirical findings regarding gate effectiveness and metric evolution. Finally, Section \ref{sec:discussion} discusses implications and threats to validity, followed by conclusions in Section \ref{sec:conclusion}.

